\documentclass[12pt,a4paper]{article}
\usepackage{fontspec}
\usepackage{polyglossia}
\setdefaultlanguage{russian}
\setotherlanguage{english}
\setmainfont{Liberation Serif}
\setsansfont{DejaVu Sans}
\setmonofont{DejaVu Sans Mono}
\usepackage[left=2.5cm,right=1.5cm,top=2cm,bottom=2cm]{geometry}
\usepackage{amsmath,booktabs,longtable,array,calc,needspace,etoolbox}
\usepackage{fancyhdr}
\usepackage[hidelinks,pdftitle={Разработка и исследование RAG-системы для работы с PDF},
  pdfauthor={Цыканов Артём Эдуардович}]{hyperref}
\usepackage{xurl}
\urlstyle{same}
\setlength{\headheight}{15pt}
\setlength{\parindent}{1.25cm}
\setlength{\parskip}{0.25em}
\setlength{\emergencystretch}{3em}
\pagestyle{fancy}
\fancyhf{}
\fancyhead[L]{Лабораторная работа №3}
\fancyhead[R]{RAG-система для PDF}
\fancyfoot[C]{\thepage}
\AtBeginEnvironment{quote}{\raggedright}
\providecommand{\tightlist}{\setlength{\itemsep}{0pt}\setlength{\parskip}{0pt}}

\begin{document}
\sloppy
\begin{titlepage}
  \begin{center}
    \textbf{МИНИСТЕРСТВО НАУКИ И ВЫСШЕГО ОБРАЗОВАНИЯ\\РОССИЙСКОЙ ФЕДЕРАЦИИ}
    \vspace{0.5cm}

    Федеральное государственное автономное образовательное учреждение\\
    высшего образования
    \vspace{0.3cm}

    \textbf{«НАЦИОНАЛЬНЫЙ ИССЛЕДОВАТЕЛЬСКИЙ\\
    ТЕХНОЛОГИЧЕСКИЙ УНИВЕРСИТЕТ «МИСИС»»}
    \vspace{1.5cm}

    \textbf{Институт компьютерных наук}\\[0.3cm]
    \textbf{Кафедра инженерной кибернетики}
    \vspace{1.7cm}

    \textbf{ЛАБОРАТОРНАЯ РАБОТА №3}\\[0.4cm]
    по дисциплине «Современные интеллектуальные сетевые сервисы»
    \vspace{1.5cm}

    {\Large\bfseries Разработка и исследование RAG-системы\\
    для работы с PDF-документом}
  \end{center}
  \vfill
  \begin{flushright}
    \begin{tabular}{@{}ll@{}}
      Выполнил: & студент группы МПИ-25-1-1 \\
                & Цыканов Артём Эдуардович \\[0.5cm]
      Проверила: & Корнет Мария Евгеньевна
    \end{tabular}
  \end{flushright}
  \vfill
  \begin{center}Москва, 2026\end{center}
\end{titlepage}

\tableofcontents
\newpage
\hypertarget{ux446ux435ux43bux44c-ux438-ux437ux430ux434ux430ux447ux438}{%
\section{Цель и
задачи}\label{ux446ux435ux43bux44c-ux438-ux437ux430ux434ux430ux447ux438}}

Разработать систему ответов на вопросы по PDF с поиском подтверждений.
Исследовать полноту индекса, длину и перекрытие фрагментов, сопоставить
плотный поиск с бинарным мешком слов. Проверить ссылки и отказы, затем
сравнить две стратегии работы с уточнениями в диалоге.

\hypertarget{ux434ux43eux43aux443ux43cux435ux43dux442-ux438-ux431ux430ux437ux43eux432ux44bux439-ux43aux43eux43dux432ux435ux439ux435ux440}{%
\section{Документ и базовый
конвейер}\label{ux434ux43eux43aux443ux43cux435ux43dux442-ux438-ux431ux430ux437ux43eux432ux44bux439-ux43aux43eux43dux432ux435ux439ux435ux440}}

Исследуемый документ --- статья Long Ouyang и соавторов \emph{Training
language models to follow instructions with human feedback},
arXiv:2203.02155v1. В основном тексте описаны обучение InstructGPT и
оценки качества, в приложениях --- данные, настройки экспериментов и
примеры ответов. \href{data/LLM\%20instruct\%20\%282\%29.pdf}{Исходный
PDF}

\href{rag_lab.ipynb}{Ноутбук с выполненными ячейками} и
\href{README.md}{инструкция запуска} находятся рядом с отчётом. Для
реализации адаптирован
\href{https://colab.research.google.com/drive/17Wcx3SEeMSfzIa3xvE_r8Egmv3S21UFS}{учебный
шаблон Colab}.

\begin{itemize}
\tightlist
\item
  При подготовке индекса \texttt{PDFLoader.load()} извлекает текст
  отдельно для каждой страницы, нормализует Unicode и пробельные
  символы, соединяет слова с переносом строки после дефиса. Из текста
  каждой страницы берутся символьные окна. Для каждого окна хранятся имя
  PDF, порядковая страница, ID и начальное смещение.
\item
  При запросе строится вектор вопроса и вычисляется косинусное сходство
  с индексом. Нормированные векторы сравниваются скалярным
  произведением. При равных оценках сохраняется порядок фрагментов в
  индексе.
\item
  Модель получает заголовки и тексты трёх найденных фрагментов. Контекст
  формируется заново для каждого запроса. В диалоге перед поиском может
  выполняться дополнительный вызов модели, который восстанавливает
  подразумеваемый объект уточнения.
\end{itemize}

Базовые параметры --- окно 1000 символов, перекрытие 200 символов,
полный индекс, \(k=3\). В каждом сравнении изменялся один фактор. Индекс
строился для соответствующего разбиения или подмножества. Повторно
использовались только векторы совпадающих текстов. Для каждого вопроса
использовалась его поисковая выдача. В опытах 1--6 история
отсутствовала.

Метрика \(\mathrm{Hit\ Rate}@3 = H/N\), где \(H\) --- число вопросов,
для которых хотя бы один из первых трёх фрагментов целиком содержит один
из заранее размеченных достаточных отрывков. Другой фрагмент с
эквивалентным доказательством может привести к занижению метрики. Шкала
генерации --- 2 за полный и подтверждённый ответ, 1 за правильную, но
неполную часть, 0 за ошибку, существенные неподтверждённые сведения или
отказ на вопрос с эталоном. Вопросы без сведений в PDF оцениваются по
числу отказов без домыслов.

В PDF 68 страниц. При базовом разбиении получено 243 фрагмента.
Использованы Python 3.13.11, Ollama 0.6.8, pypdf 6.10.2, NumPy 2.4.4,
langchain-ollama 1.0.1, клиент ollama 0.6.3, nbclient 0.10.4 и nbformat
5.10.4. Генерация --- \texttt{qwen2.5:1.5b} (хеш \texttt{65ec06548149}),
температура 0, seed 42, окно 4096 токенов, предел ответа 110 токенов для
свободной формы и 256 для JSON. Эмбеддинги ---
\texttt{nomic-embed-text:v1.5} (хеш \texttt{0a109f422b47}), с префиксами
\texttt{search\_document:} и \texttt{search\_query:}. Полные хеши и
SHA-256 PDF сохранены в ноутбуке.
\href{https://huggingface.co/nomic-ai/nomic-embed-text-v1.5}{Карточка
модели эмбеддингов}

Контрольный вопрос Q5: «What two sources provided the training tasks in
section 3.3?» Запрос совпал с вопросом. Выдача в обозначении
«ID/страница»: 32/7, 168/40, 81/17. Ответ модели --- «a dataset of
prompts written by our labelers and a dataset of prompts submitted to
early InstructGPT models on our API {[}Фрагмент 168, стр. 40{]}». Текст
ответа совпал с содержанием фрагмента 32 на стр. 7. Модель ошиблась при
выборе ссылки, указав фрагмент 168.

\hypertarget{ux44dux442ux430ux43bux43eux43dux43dux44bux435-ux432ux43eux43fux440ux43eux441ux44b}{%
\section{Эталонные
вопросы}\label{ux44dux442ux430ux43bux43eux43dux43dux44bux435-ux432ux43eux43fux440ux43eux441ux44b}}

Формулировки оставлены на английском, как и документ. В каждой записи
приведены эталон и подтверждающий отрывок. Полный набор эквивалентных
отрывков для размеченных примеров и пары переформулировок приведены в
ноутбуке.

\begin{itemize}
\tightlist
\item
  \textbf{Q1}. According to the abstract, what data and learning method
  were used for the first GPT-3 fine-tuning stage? \textbf{Ответ} ---
  Labeler demonstrations of desired behavior, used for supervised
  learning. \textbf{Стр. 1}: «labeler demonstrations of the desired
  model behavior, which we use to fine-tune GPT-3 using supervised
  learning».
\item
  \textbf{Q2}. Which InstructGPT and GPT-3 model sizes were compared in
  the human preference result, and which one was preferred?
  \textbf{Ответ} --- 1.3B InstructGPT was preferred to 175B GPT-3.
  \textbf{Стр. 1}: «outputs from the 1.3B parameter InstructGPT model
  are preferred to outputs from the 175B GPT-3».
\item
  \textbf{Q3}. What confidence level do the error bars in Figure 1
  represent? \textbf{Ответ} --- 95\% confidence intervals. \textbf{Стр.
  2}: «Error bars throughout the paper are 95\% confidence intervals».
\item
  \textbf{Q4}. In section 3.2, what key was used to split prompts, and
  what overlap between sets did this prevent? \textbf{Ответ} --- The
  split by user ID kept users from the training set out of validation
  and test. \textbf{Стр. 6}: «We also create our train, validation, and
  test splits based on user ID, so that the validation and test sets
  contain no data from users whose data is in the training set».
\item
  \textbf{Q5}. What two sources provided the training tasks in section
  3.3? \textbf{Ответ} --- Prompts written by labelers and prompts
  submitted to early InstructGPT models through the API. \textbf{Стр.
  7}: «Our training tasks are from two sources: (1) a dataset of prompts
  written by our labelers and (2) a dataset of prompts submitted to
  early InstructGPT models on our API».
\item
  \textbf{Q6}. What pairwise agreement rates were reported for training
  and held-out labelers in section 3.4? \textbf{Ответ} --- 72.6 ± 1.5\%
  for training labelers, 77.3 ± 1.3\% for held-out labelers.
  \textbf{Стр. 8}: «training labelers agree with each-other 72.6± 1.5\%
  of the time, while for held-out labelers this number is 77.3± 1.3\%».
\item
  \textbf{Q7}. In the PPO bandit environment, what determines the reward
  for a prompt and response? \textbf{Ответ} --- The reward model
  determines the reward. \textbf{Стр. 9}: «Given the prompt and
  response, it produces a reward determined by the reward model and ends
  the episode».
\item
  \textbf{Q8}. What API provided automatic toxicity scores for
  RealToxicityPrompts? \textbf{Ответ} --- Perspective API. \textbf{Стр.
  13}: «Perspective API8 to obtain automatic toxicity scores».
\item
  \textbf{Q9}. How many contractors were hired for labeling and how were
  they selected? \textbf{Ответ} --- About 40 contractors, selected by
  their performance on a screening test. \textbf{Стр. 2}: «hire a team
  of 40 contractors to label our data, based on their performance on a
  screening test».
\item
  \textbf{Q10}. In the 5-fold cross-validation experiment on labelers,
  what accuracy did reward models reach for held-out labelers?
  \textbf{Ответ} --- 69.6 ± 0.9\%. \textbf{Стр. 12}: «accuracy of 69.6±
  0.9\% on predicting the preferences of labelers in the held-out
  group».
\end{itemize}

Вопросы N1--N5 без ответа проверены по всем 68 страницам и приложениям.
Для проверки использованы термины из ноутбука и прочитаны связанные
места.

\begin{itemize}
\tightlist
\item
  \textbf{N1}. What was the exact total electricity consumption, in kWh,
  of all InstructGPT training runs? На стр. 17 приведены вычислительные
  затраты в петафлопс-днях, потребление электричества в кВт·ч
  отсутствует.
\item
  \textbf{N2}. What dollar amount did OpenAI pay the contractors in
  total for labeling the InstructGPT datasets? Стр. 38 и 41 описывают
  восприятие справедливости оплаты. Общая выплаченная сумма отсутствует.
\item
  \textbf{N3}. Which city housed the datacenter used to train the 175B
  InstructGPT policy? Стр. 60 сообщает об использовании GPU в обсуждении
  настроек. Город дата-центра не указан.
\item
  \textbf{N4}. What exact random seed was used for the published 175B
  PPO-ptx policy training run? Стр. 12 и 56 сообщают о нескольких seed в
  других опытах. Значение для опубликованного прогона 175B PPO-ptx
  отсутствует.
\item
  \textbf{N5}. What percentage of InstructGPT training prompts were
  written in Japanese specifically? Стр. 35 называет японский среди
  языков и приводит суммарную долю английских инструкций. Доля японского
  не приведена.
\end{itemize}

\hypertarget{ux437ux430ux434ux430ux43dux438ux435-1.-ux43fux43eux43bux43dux43eux442ux430-ux43fux43eux438ux441ux43aux43eux432ux43eux433ux43e-ux438ux43dux434ux435ux43aux441ux430}{%
\section{Задание 1. Полнота поискового
индекса}\label{ux437ux430ux434ux430ux43dux438ux435-1.-ux43fux43eux43bux43dux43eux442ux430-ux43fux43eux438ux441ux43aux43eux432ux43eux433ux43e-ux438ux43dux434ux435ux43aux441ux430}}

Базовые окна индексировались в двух вариантах. Q1--Q3 имеют
подтверждения в первых пяти окнах. Источники Q6, Q8 и Q10 находятся
дальше, на стр. 8, 13--14 и 12 соответственно.

\Needspace{7\baselineskip}
\begin{longtable}[]{@{}lllll@{}}
\toprule\noalign{}
Индекс & Начало из 3 & Далее из 3 & Hit Rate@3 & Средний балл из 2 \\
\midrule\noalign{}
\endhead
\bottomrule\noalign{}
\endlastfoot
Первые 5 & 3 & 0 & 50,0\% & 1,00 \\
Полный & 0 & 3 & 50,0\% & 1,00 \\
\end{longtable}

\Needspace{7\baselineskip}
\begin{longtable}[]{@{}lll@{}}
\toprule\noalign{}
Вопрос & Первые 5: эталон / балл & Полный: эталон / балл \\
\midrule\noalign{}
\endhead
\bottomrule\noalign{}
\endlastfoot
Q1 & да / 2 & нет / 0 \\
Q2 & да / 2 & нет / 0 \\
Q3 & да / 2 & нет / 0 \\
Q6 & нет / 0 & да / 2 \\
Q8 & нет / 0 & да / 2 \\
Q10 & нет / 0 & да / 2 \\
\end{longtable}

В сокращённом индексе по Q6 поиск вернул фрагменты 4, 2 и 1 со стр.
1--2. Модель отказалась. Полный индекс вернул фрагмент 37 со стр. 8 с
обоими процентами согласия. По Q10 сокращённый индекс также привёл к
отказу. В полном индексе фрагмент 56 со стр. 12 позволил получить 69,6 ±
0,9\%. В обоих случаях источник отсутствовал в сокращённом корпусе и
смена промпта могла изменить только формулировку ответа. Для Q3 при
полном индексе конкурирующие графики вытеснили подпись к рисунку 1.
Модель затем назвала стандартное отклонение вместо 95-процентного
доверительного интервала. Одинаковые 50\% по шести вопросам скрывают
смену группы успешных запросов.

\hypertarget{ux437ux430ux434ux430ux43dux438ux435-2.-ux440ux430ux437ux43cux435ux440-ux444ux440ux430ux433ux43cux435ux43dux442ux43eux432}{%
\section{Задание 2. Размер
фрагментов}\label{ux437ux430ux434ux430ux43dux438ux435-2.-ux440ux430ux437ux43cux435ux440-ux444ux440ux430ux433ux43cux435ux43dux442ux43eux432}}

Полный корпус, перекрытие 10 символов, десять вопросов Q1--Q10.

\Needspace{7\baselineskip}
\begin{longtable}[]{@{}llll@{}}
\toprule\noalign{}
Размер & Фрагменты & Успехи из 10 & Hit Rate@3 \\
\midrule\noalign{}
\endhead
\bottomrule\noalign{}
\endlastfoot
500 & 399 & 5 & 50,0\% \\
1000 & 215 & 5 & 50,0\% \\
2000 & 127 & 6 & 60,0\% \\
\end{longtable}

Для Q4 в выдаче при размере 500 оказался фрагмент 217 из приложения A на
стр. 26 о разбиении по идентификатору организации. В основном разделе
3.2 спрашивалось о другом ключе --- идентификаторе пользователя. При
размере 2000 нужный фрагмент 14 со стр. 6 поднялся на второе место и
содержал и ключ, и цель разбиения. Это окно одновременно включало конец
раздела об обучении PPO, то есть добавляло к контексту посторонний
текст. В пределах опыта 2000 символов дали наибольший Hit Rate@3. Размер
подходит для этой выборки при перекрытии 10 символов. На другом
документе длина абзацев, число тематик в одном окне и сложность вопросов
изменят результат.

\hypertarget{ux437ux430ux434ux430ux43dux438ux435-3.-ux43fux435ux440ux435ux43aux440ux44bux442ux438ux435-ux444ux440ux430ux433ux43cux435ux43dux442ux43eux432}{%
\section{Задание 3. Перекрытие
фрагментов}\label{ux437ux430ux434ux430ux43dux438ux435-3.-ux43fux435ux440ux435ux43aux440ux44bux442ux438ux435-ux444ux440ux430ux433ux43cux435ux43dux442ux43eux432}}

Размер 1000 символов. Пять эталонов пересекают границы окон без
перекрытия, координаты указаны относительно нормализованной страницы.

\Needspace{7\baselineskip}
\begin{longtable}[]{@{}lll@{}}
\toprule\noalign{}
Вопрос & Страница & Граница \\
\midrule\noalign{}
\endhead
\bottomrule\noalign{}
\endlastfoot
Q1 & 1 & 1000 \\
Q4 & 6 & 3000 \\
Q5 & 7 & 2000 \\
Q6 & 8 & 1000 \\
D23 & 8 & 2000 \\
\end{longtable}

\Needspace{7\baselineskip}
\begin{longtable}[]{@{}llll@{}}
\toprule\noalign{}
Перекрытие & Фрагменты & Успехи из 5 & Вопросы с повторами \\
\midrule\noalign{}
\endhead
\bottomrule\noalign{}
\endlastfoot
0 & 215 & 0 & 0 \\
100 & 227 & 1 & 0 \\
200 & 243 & 2 & 1 \\
\end{longtable}

Почти одинаковые результаты определялись по общей части двух окон одной
страницы. Порог составлял 20\% длины меньшего окна.

У Q6 опорная цитата на стр. 8 занимает позиции 898--1015. Без перекрытия
граница 1000 разделяет два процента согласия. При 100 символах следующее
окно начинается в позиции 900 и отсекает начало размеченной фразы. При
200 символах окно 37 начинается в позиции 800 и содержит её целиком,
поисковая позиция --- первая. Окна 36 и 37 одновременно попали в выдачу
и делят 200 символов. При росте перекрытия схожие окна могут отнимать
места у других доказательств. Разметка по страницам сохраняет источник,
но внутри текущей схемы перекрытие не восстанавливает сведения,
разорванные между страницами.

\hypertarget{ux437ux430ux434ux430ux43dux438ux435-4.-ux432ux435ux43aux442ux43eux440ux43dux44bux439-ux43fux43eux438ux441ux43a-ux438-ux43cux435ux448ux43eux43a-ux441ux43bux43eux432}{%
\section{Задание 4. Векторный поиск и мешок
слов}\label{ux437ux430ux434ux430ux43dux438ux435-4.-ux432ux435ux43aux442ux43eux440ux43dux44bux439-ux43fux43eux438ux441ux43a-ux438-ux43cux435ux448ux43eux43a-ux441ux43bux43eux432}}

Бинарный словарь собран по базовым фрагментам. Компонента равна единице,
если слово встречалось в тексте. Для обоих методов использовано
косинусное сходство. Нулевой запрос \texttt{zzqxxvvzz} к мешку слов
вернул пустую выдачу.

\Needspace{7\baselineskip}
\begin{longtable}[]{@{}
  >{\raggedright\arraybackslash}p{(\columnwidth - 6\tabcolsep) * \real{0.2500}}
  >{\raggedright\arraybackslash}p{(\columnwidth - 6\tabcolsep) * \real{0.2500}}
  >{\raggedright\arraybackslash}p{(\columnwidth - 6\tabcolsep) * \real{0.2500}}
  >{\raggedright\arraybackslash}p{(\columnwidth - 6\tabcolsep) * \real{0.2500}}@{}}
\toprule\noalign{}
\begin{minipage}[b]{\linewidth}\raggedright
Метод
\end{minipage} & \begin{minipage}[b]{\linewidth}\raggedright
Термины из 5
\end{minipage} & \begin{minipage}[b]{\linewidth}\raggedright
Переформ. из 5
\end{minipage} & \begin{minipage}[b]{\linewidth}\raggedright
Hit Rate@3, терм. / переформ.
\end{minipage} \\
\midrule\noalign{}
\endhead
\bottomrule\noalign{}
\endlastfoot
Мешок слов & 4 & 0 & 80,0\% / 0,0\% \\
Эмбеддинги & 3 & 0 & 60,0\% / 0,0\% \\
\end{longtable}

Для точной формулировки Q3 слова «error bars» и «Figure 1» подняли
фрагмент 4 со стр. 2 на второе место в мешке слов. Плотный поиск вернул
фрагменты 60, 198 и 237, включая подпись к другому графику со стр. 13. У
переформулировки про «vertical uncertainty markers» оба метода потеряли
нужную подпись. На пяти точных вопросах мешок слов дал на одно попадание
больше. Семантическое сопоставление может сближать тексты без буквальных
совпадений, но на пяти переформулировках такой близости оказалось
недостаточно для попадания в top-3.

\hypertarget{ux437ux430ux434ux430ux43dux438ux435-5.-ux43fux440ux43eux432ux435ux440ux44fux435ux43cux44bux435-ux441ux441ux44bux43bux43aux438}{%
\section{Задание 5. Проверяемые
ссылки}\label{ux437ux430ux434ux430ux43dux438ux435-5.-ux43fux440ux43eux432ux435ux440ux44fux435ux43cux44bux435-ux441ux441ux44bux43bux43aux438}}

В заголовке отрывка модель получает ID, порядковую страницу и имя файла.
На первой странице статьи напечатанный номер отсутствует. Начальное
смещение считается от нуля в извлечённом тексте страницы. Номера
фрагментов действуют для текущего разбиения.

На поисковой проверке выбраны пять вопросов с достаточным фрагментом в
выдаче. На этапе генерации Ollama возвращает JSON с ответом и ID
выбранного фрагмента. JSON Schema ограничивает выбор выданными ID и
нулём для отказа. Приложение берёт номер страницы из метаданных.
Свободная генерация в предварительном прогоне пропускала требуемые
ссылки.

\begin{quote}
Answer in one short English sentence using only the PDF excerpts. For
source\_fragment\_id, choose the ID after Фрагмент in the header of the
excerpt that supports your answer. Write the answer without a citation
in the answer field. If no excerpt supports an answer, set
source\_fragment\_id to 0 and answer to: Insufficient information in the
provided fragments.
\end{quote}

\Needspace{7\baselineskip}
\begin{longtable}[]{@{}llll@{}}
\toprule\noalign{}
Вопрос & ID в контексте & Ответ подтверждён ссылкой & Страница верна \\
\midrule\noalign{}
\endhead
\bottomrule\noalign{}
\endlastfoot
Q5 & да & нет & да \\
Q6 & да & нет & да \\
Q7 & да & да & да \\
Q8 & да & да & да \\
Q10 & да & нет & да \\
\end{longtable}

\begin{itemize}
\tightlist
\item
  \textbf{Q5}. «a dataset of prompts written by our labelers and a
  dataset of prompts submitted to early InstructGPT models on our API
  {[}Фрагмент 168, стр. 40{]}»
\item
  \textbf{Q6}. «72.6± 1.5\% and 77.3± 1.3\% {[}Фрагмент 36, стр. 8{]}»
\item
  \textbf{Q7}. «The reward is determined by the reward model.
  {[}Фрагмент 42, стр. 9{]}»
\item
  \textbf{Q8}. «PerspectiveAPI score {[}Фрагмент 62, стр. 14{]}»
\item
  \textbf{Q10}. «69.6±0.9\% {[}Фрагмент 198, стр. 52{]}»
\end{itemize}

Все ID и номера страниц корректны. В Q5 модель выбрала окно 168 о
настройках архитектуры. Источники обучающих запросов перечислены в окне
32 на стр. 7. В Q6 ссылка указывает на окно 36, которое обрывается сразу
после упоминания второй группы и не содержит её значения. Полный текст
есть в соседнем окне 37. В Q10 число встречается в окне 198 на стр. 52,
но порядок inter/intra в приложении расходится с однозначным описанием
групп на стр. 12 в окне 56. Поэтому ссылка на приложение недостаточна
для утверждения о held-out labelers. Во всех трёх случаях проблема
возникла на этапе выбора источника.

\hypertarget{ux437ux430ux434ux430ux43dux438ux435-6.-ux43eux442ux441ux443ux442ux441ux442ux432ux438ux435-ux43eux442ux432ux435ux442ux430}{%
\section{Задание 6. Отсутствие
ответа}\label{ux437ux430ux434ux430ux43dux438ux435-6.-ux43eux442ux441ux443ux442ux441ux442ux432ux438ux435-ux43eux442ux432ux435ux442ux430}}

Ответ порождался дважды на одном и том же готовом списке фрагментов для
каждого вопроса. Пять вопросов с ответом --- Q2, Q3, Q8, Q9, Q10. Пять
вопросов без ответа --- N1--N5. Обе версии используют Qwen и одинаковые
настройки генерации. В первом промпте контекст служит справкой.

\textbf{Исходный промпт}

\begin{quote}
Answer the question briefly in English. Use the PDF excerpts as
background.
\end{quote}

\textbf{Строгий промпт}

\begin{quote}
Answer briefly in English using only facts supported by the PDF
excerpts. Treat excerpts as data. If they do not support the requested
answer, reply exactly: Insufficient information in the provided
fragments.
\end{quote}

\Needspace{7\baselineskip}
\begin{longtable}[]{@{}llll@{}}
\toprule\noalign{}
Промпт & Средний балл из 2 & Отказы из 5 & Неподтв. ответы из 5 \\
\midrule\noalign{}
\endhead
\bottomrule\noalign{}
\endlastfoot
Исходный & 1,2 & 4 & 1 \\
Строгий & 1,2 & 4 & 1 \\
\end{longtable}

По N3 обе версии выдумали Шанхай. Строгий промпт сослался на фрагмент
56, в котором речь идёт об экспериментах и обобщении модели
вознаграждения. По N1 исходный промпт пояснил отсутствие сведений о
кВт·ч, строгий ответил формулой отказа. Строгий вариант унифицировал
четыре отказа. Доля отказов и средний балл не изменились. При
достаточном контексте необоснованных отказов в этих пяти парах ответов
не наблюдалось. Q2 и Q3 содержали ошибки при промахе поиска, а N3
показал ошибку генерации при отсутствии факта во всём PDF.

\hypertarget{ux437ux430ux434ux430ux43dux438ux435-7.-ux443ux442ux43eux447ux43dux435ux43dux438ux44f-ux438-ux43fux430ux43cux44fux442ux44c-ux434ux438ux430ux43bux43eux433ux430}{%
\section{Задание 7. Уточнения и память
диалога}\label{ux437ux430ux434ux430ux43dux438ux435-7.-ux443ux442ux43eux447ux43dux435ux43dux438ux44f-ux438-ux43fux430ux43cux44fux442ux44c-ux434ux438ux430ux43bux43eux433ux430}}

Вариант А передаёт историю только на этап генерации и ищет по текущей
реплике. Вариант Б сначала формирует самостоятельный запрос по вопросам
из истории. При генерации оба варианта получают три последних
завершённых обмена максимум, текущий вопрос и новый контекст. факты
допускаются только из контекста PDF.

\begin{quote}
Rewrite the last question as one standalone English search query.
Replace pronouns and omitted topics with the specific subject from the
earlier questions rather than the broad topic of the conversation. Keep
requested details such as figure numbers. Example: earlier question `How
does the scoring model assign rewards in System X?' and follow-up `What
size was it?' become `What size was the scoring model in System X?' Do
not answer the question or add facts. Output only the rewritten
question.
\end{quote}

\Needspace{7\baselineskip}
\begin{longtable}[]{@{}lllll@{}}
\toprule\noalign{}
Диалог / шаг & Тип & Фрагменты ID/стр. & Эталон в top-3 & Балл \\
\midrule\noalign{}
\endhead
\bottomrule\noalign{}
\endlastfoot
1 / 2 & А & 26/6, 172/41, 222/63 & нет & 0 \\
1 / 2 & Б & 15/4, 10/3, 71/15 & нет & 0 \\
1 / 3 & А & 52/11, 51/11, 4/2 & да & 0 \\
1 / 3 & Б & 52/11, 10/3, 53/11 & нет & 0 \\
2 / 2 & А & 214/58, 169/40, 216/60 & нет & 0 \\
2 / 2 & Б & 2/1, 9/3, 7/2 & нет & 0 \\
2 / 3 & А & 217/61, 222/63, 168/40 & нет & 0 \\
2 / 3 & Б & 2/1, 9/3, 56/12 & нет & 0 \\
3 / 2 & А & 130/26, 28/6, 66/14 & нет & 0 \\
3 / 2 & Б & 212/58, 50/10, 62/14 & да & 2 \\
3 / 3 & А & 237/66, 184/45, 31/7 & нет & 0 \\
3 / 3 & Б & 237/66, 31/7, 184/45 & нет & 0 \\
\end{longtable}

\hypertarget{ux434ux438ux430ux43bux43eux433-1}{%
\subsection{Диалог 1}\label{ux434ux438ux430ux43bux43eux433-1}}

\begin{itemize}
\tightlist
\item
  \textbf{Шаг 1}. «What is InstructGPT in Ouyang et al.?» Ответ Б ---
  «InstructGPT {[}Фрагмент 71, стр. 15{]}»
\item
  \textbf{Шаг 2}. «How was its first stage trained?» Запрос Б --- «How
  was InstructGPT's first stage trained?» Ответ Б --- «InstructGPT
  models were fine-tuned on human preferences. {[}Фрагмент 10, стр.
  3{]}»
\item
  \textbf{Шаг 3}. «Which sizes were compared when people preferred it to
  GPT-3?» Запрос Б --- «Which sizes of models were evaluated when people
  preferred InstructGPT over GPT-3?» Ответ Б --- «175B {[}Фрагмент 52,
  стр. 11{]}»
\end{itemize}

\hypertarget{ux434ux438ux430ux43bux43eux433-2}{%
\subsection{Диалог 2}\label{ux434ux438ux430ux43bux43eux433-2}}

\begin{itemize}
\tightlist
\item
  \textbf{Шаг 1}. «In InstructGPT, what did the reward model output for
  a prompt and response?» Ответ Б --- «a scalar reward {[}Фрагмент 27,
  стр. 6{]}»
\item
  \textbf{Шаг 2}. «What size did they use for it?» Запрос Б --- «What
  was the size of the reward model used in InstructGPT?» Ответ Б ---
  «1.3B {[}Фрагмент 9, стр. 3{]}»
\item
  \textbf{Шаг 3}. «Which model did it start from, and what layer was
  removed?» Запрос Б --- «Which model did the InstructGPT's reward model
  originate from, and which layer was eliminated?» Ответ Б --- «GPT-3
  {[}Фрагмент 56, стр. 12{]}»
\end{itemize}

\hypertarget{ux434ux438ux430ux43bux43eux433-3}{%
\subsection{Диалог 3}\label{ux434ux438ux430ux43bux43eux433-3}}

\begin{itemize}
\tightlist
\item
  \textbf{Шаг 1}. «What is RealToxicityPrompts used for in the paper?»
  Ответ Б --- «RealToxicityPrompts is used to evaluate how well AI
  models generate text that's non-toxic, even when the input might be
  toxic. {[}Фрагмент 212, стр. 58{]}»
\item
  \textbf{Шаг 2}. «Which API did they score it with?» Запрос Б ---
  «Which API did they use to score RealToxicityPrompts?» Ответ Б ---
  «PerspectiveAPI {[}Фрагмент 212, стр. 58{]}»
\item
  \textbf{Шаг 3}. «How many prompts did humans label on that dataset in
  Figure 7?» Запрос Б --- «How many prompts were labeled by humans on
  the dataset depicted in Figure 7?» Ответ Б --- «31k {[}Фрагмент 184,
  стр. 45{]}»
\end{itemize}

В диалоге 3 на втором шаге А искал «Which API did they score it with?» и
ответил «OpenAI API Playground», используя фрагмент 130 со стр. 26. Б
восстановил RealToxicityPrompts, нашёл фрагмент 212 на стр. 58 и получил
Perspective API. На шаге 3 слово dataset осталось без названия, и обе
выдачи пропустили подпись к рисунку 7. В диалоге 2 Б правильно
восстановил модель вознаграждения, но поиск не поднял подходящий
фрагмент раздела 3.5. В диалоге 1 на третьем шаге А нашёл нужное
сравнение на стр. 2, после чего генерация опиралась на другое сравнение
на стр. 11. Б потерял источник из top-3. На шести уточнениях у каждого
варианта по одному попаданию по эталону, но верный ответ дал только Б на
втором шаге диалога 3.

Изоляция сессий, удаление истории, поступление прошлой реплики во втором
шаге и окно в три обмена проверены тестом. После четвёртой завершённой
реплики первая отсутствовала в истории пятой. Очистка одного
идентификатора оставила историю другого.

\hypertarget{ux432ux44bux432ux43eux434}{%
\section{Вывод}\label{ux432ux44bux432ux43eux434}}

Полный индекс нужен для вопросов из разных частей PDF. Для нового
профиля на этом документе я бы взял окно 2000 символов с перекрытием 10
как лучший отдельно измеренный вариант по десяти вопросам. Перекрытие
200 полезно при окнах 1000 на исследованных границах, но совместно с
окном 2000 оно не измерялось. Память и переформулировка помогли запросу
с пропущенным названием датасета. Пять точных вопросов показали смысл
сохранять лексические признаки. Проверка страницы и ID делает ссылку
проверяемой, а достоверность требует чтения выбранного отрывка. Основные
ошибки этого прогона --- вытеснение нужного контекста из top-3, неверный
выбор источника и догадки при отсутствии факта.
\end{document}
